

\subsection*{Punto2: Implementación de la Modulación M-PSK}
Para implementar la modulación M-PSK empleando los métodos VCO y tablas de verdad, se configuró el flujo de procesamiento en GNU Radio como se muestra en la figura 1. Se utilizó una frecuencia de muestreo de 85.3333 kHz y un factor de 8 muestras por símbolo ($\text{sps}=8$). La tasa de símbolos del sistema se define teóricamente como:
$$R_s=\frac{R_b}{\log_2(M)}=\frac{32000}{3}=10666.67\text{ baudios}$$
\begin{figure}[H]
    \centering
    \includegraphics[width=0.9\linewidth]{/home/valentina/ComII_G2_equipo2/practica_4/Informe/imagenes/1.jpeg}
    \caption{flujograma}
\end{figure}
La modulación se generó mediante dos enfoques paralelos para evidenciar su equivalencia en el desarrollo de la práctica:
\begin{enumerate}
    \item \textbf{Método VCO (Oscilador Controlado por Voltaje):} La secuencia discreta de símbolos $m\in\texttt{Random Source}$ ingresa al bloque \texttt{VCO complex}, el cual varía la fase de la señal portadora para representar la información. La fase teórica asignada sigue la relación:
    $$\theta_m=\frac{2\pi m}{8}$$
    \item \textbf{Método de Tabla de Verdad (Mapeo Directo):} Tal como se indica en la guía, la tabla de verdad en GNU Radio corresponde a los puntos de la constelación escritos en forma de un vector. Se utilizó el bloque \texttt{Chunks to Symbols} para asignar a cada símbolo entero su valor cartesiano en el plano complejo. El ordenamiento de los símbolos coincide con el vector de la constelación, tal como se detalla en el Cuadro \ref{tab:tabla_verdad}.
\end{enumerate}
\begin{table}[h!]
\centering
\caption{Tabla de verdad y coordenadas de la constelación 8-PSK.}
\label{tab:tabla_verdad}
\begin{tabular}{|c|c|c|c|}
\hline
\textbf{Bits} & \textbf{Símbolo ($m$)} & \textbf{Fase ($\theta_m$)} & \textbf{Coordenada Compleja ($I+jQ$)} \\ \hline
000 & 0 & $0^\circ$ & $1+0j$ \\ \hline
001 & 1 & $45^\circ$ & $0.707+0.707j$ \\ \hline
010 & 2 & $90^\circ$ & $0+1j$ \\ \hline
011 & 3 & $135^\circ$ & $-0.707+0.707j$ \\ \hline
100 & 4 & $180^\circ$ & $-1+0j$ \\ \hline
101 & 5 & $225^\circ$ & $-0.707-0.707j$ \\ \hline
110 & 6 & $270^\circ$ & $0-1j$ \\ \hline
111 & 7 & $315^\circ$ & $0.707-0.707j$ \\ \hline
\end{tabular}
\end{table}
\subsection*{Punto 3: Ancho de Banda de la Envolvente Compleja}
Para obtener el ancho de banda de la envolvente compleja, se analizó el espectro de la señal en banda base. La medición teórica y práctica se enfoca en la distancia desde la frecuencia central hasta el primer nulo de potencia.
\subsection*{Punto 4: Señal, PSD y Cruces por Cero}

Se empleó el sumidero de frecuencia \texttt{QT GUI Frequency Sink} para mostrar la señal y su PSD. Se procedió a identificar la frecuencia central y a registrar los valores exactos en los que el espectro pasa por cero.


% Recuerda incluir aquí tu imagen con: \includegraphics[width=0.8\textwidth]{nombre_de_captura_psd.png}

